\section{Coaccessibility-Aware Sharding}
\label{sec:offline_sharding}

This section formalizes the sharding objective and details how \sysname constructs its distributed index to minimize fan-out while balancing system load.

\subsection{Formulation: Hypergraph Min-Cut}
Let $X=\{x_1,\dots,x_N\}$ be the dataset, $q\sim Q$ a query drawn from the workload, and $\mathrm{NN}_k(q)\subseteq X$ its true $k$ nearest neighbors under the target metric. A strict sharding is an assignment $\sigma:X\to\{1,\dots,P\}$. To return the true neighbors, a query must touch every shard that holds at least one of them. Thus, its required fan-out is:
\[
\mathrm{fanout}^\star_\sigma(q)=\bigl|\{\sigma(x):x\in\mathrm{NN}_k(q)\}\bigr|.
\]

As established in \S\ref{sec:motivation}, fan-out is the first-order term determining aggregate system work. Our primary objective is to minimize the expected fan-out across the query workload, subject to a strict size balance constraint to distribute storage evenly:
\begin{equation}
\label{eq:objective}
\begin{aligned}
\min_{\sigma}\quad & \mathbb{E}_{q}\!\left[\mathrm{fanout}^\star_\sigma(q)\right] \\
\text{s.t.}\quad & (1-\varepsilon)\tfrac{N}{P}\le |\sigma^{-1}(i)| \le (1+\varepsilon)\tfrac{N}{P}\quad \forall i \in \{1,\dots,P\}.
\end{aligned}
\end{equation}

This objective naturally maps to a hypergraph $H=(X,\{e_q\}_q)$. The vertices are the vector set $X$, and each query's target set forms a hyperedge $e_q=\mathrm{NN}_k(q)$. Intuitively, if a query's target neighbors are scattered across many different shards, its hyperedge is ``cut'' multiple times, forcing the system to fan out to all those shards to retrieve them. Conversely, if all target neighbors are co-located in a single shard, the hyperedge remains uncut and the fan-out is optimal (i.e., 1). Consequently, minimizing the expected fan-out is mathematically equivalent to minimizing the standard connectivity metric (cut) in hypergraph partitioning.

\subsection{Tractable Surrogate via Navigation Graph}
\label{sec:nav_graph}
Optimizing the ideal hypergraph objective directly is impossible in practice. \emph{Offline}, future queries are unknown, so the exact hyperedges $e_q$ do not exist. Even if they did, the problem is NP-hard and inapproximable~\cite{andreev2006}.

\sysname sidesteps both obstacles by replacing the query hypergraph with a small, query-independent \emph{navigation graph} $G=(V,E,w)$ that is built entirely from the data itself. Its construction proceeds in three stages---\emph{sample}, \emph{build}, and \emph{color}.

\noindent\textbf{Sample.} Building and partitioning a graph over all $N$ points is prohibitive at billion scale. \sysname instead draws a uniform random subset $V\subseteq X$ of \emph{navigation nodes} (by default $|V|\approx N/32$, i.e.\ $\sim\!3\%$ of the data). This subset acts as a sparse skeleton that captures the global topology of the dataset while keeping the partitioning problem small.

\noindent\textbf{Build.} \sysname builds a proximity graph (an HNSW index) over $V$, connecting each node to its nearest navigation neighbors. To estimate how strongly two nodes are \emph{co-accessed} by search, \sysname uses the data points themselves as proxy queries---valid under an I.I.D. workload: each point is routed through $G$ and \emph{attached} to the navigation nodes it reaches. Aggregating these attachments defines the \emph{navigation traffic} on a directed edge $(u,v)$:
\begin{equation}
\label{eq:traffic}
w(u,v)=\mathbb{E}_{q\sim X}\bigl[\,\#\{\text{times $q$'s traversal explores }(u,v)\}\,\bigr],
\end{equation}
the expected number of times a search evaluates $v$ from $u$. The working hypothesis (\emph{search-induced topology}) is that $w$ is an accurate, computable proxy for co-relevance: points reached together during navigation tend to fall in a common query's target set. Crucially, $w$ is derived from search behavior, not from raw geometric distance.

Concretely, \sysname realizes $w$ through \emph{attachment}: for every data point it records the small set of navigation nodes its self-search reaches (its \emph{attachment set}), so an edge $(u,v)$ accumulates weight whenever $u$ and $v$ are co-attached by the same proxy query. This is the only signal the partitioner consumes---no query log or ground-truth neighbor is ever required---and it is what the offline objective below actually minimizes as the \emph{attachment cut}.

\noindent\textbf{Color.} Finally, \sysname \emph{colors} the navigation graph, assigning each node $v\in V$ a shard label $\pi(v)\in\{1,\dots,P\}$. Under this surrogate, the intractable hypergraph min-cut reduces to a standard \emph{weighted graph edge-cut}: minimize $\sum_{(u,v):\pi(u)\neq\pi(v)} w(u,v)$ subject to the balance band of Eq.~\eqref{eq:objective}. \sysname solves it with a lightweight two-phase procedure. First, a \emph{standard geometric $K$-means} over $V$ produces a locality-preserving initial coloring; the seed only needs to supply geometric coherence, since---as we show in \S\ref{sec:eval}---the deliverable fan-out is insensitive to how (or whether) the seed itself is balanced. Second, an iterative \emph{majority-vote label propagation} refines it: each node moves to the shard holding the majority of its attachment neighbors, directly lowering the attachment cut, under a capacity guard that keeps every shard within the balance band. Balance is thus owned by the refinement, not the seed, which lets \sysname adopt the simplest possible initializer.

The following lemma justifies the reduction: the weighted cut that \sysname minimizes upper-bounds the online fan-out it actually pays.

\begin{lemma}[Realized fan-out is dominated by the weighted cut]
\label{lem:fanout}
Let $\mathrm{cut}(\pi)=\{(u,v)\in E:\pi(u)\neq\pi(v)\}$ be the edges cut by the coloring $\pi$. If a query's entry points are selected from its traversal path on $G$, the realized fan-out is bounded by:
\[
\mathbb{E}_{q}\!\left[\mathrm{fanout}^\pi(q)\right]
\ \le\ 1+\!\!\sum_{(u,v)\in\mathrm{cut}(\pi)}\!\! w(u,v).
\]
\end{lemma}

\begin{proof}
For a query $q$, the set of expanded nodes $\mathrm{Vis}(q)$ and explored edges $T(q)$ form a connected component containing a spanning tree rooted at the entry node. Contracting each shard's portion of $\mathrm{Vis}(q)$ into a supernode yields a connected multigraph of $m$ supernodes (where $m$ is the number of shards visited). This multigraph must contain at least $m-1$ edges, all of which are cut edges in $T(q)$. Thus, $m \le 1 + |\{(u,v)\in T(q):\pi(u)\neq\pi(v)\}|$. Since the realized fan-out is at most $m$, taking expectations over $q$ yields the bound.
\end{proof}

\subsection{Sharding the Full Dataset}
The coloring $\pi$ partitions only the navigation skeleton $V$ ($\sim\!3\%$ of the data). \sysname propagates it to the full dataset through the same attachment signal: each base-layer (L0) point inherits the shard labels of the navigation nodes it attaches to and is assigned to the majority shard. This extends the low-cut partition from the sparse skeleton to all $N$ points at near-linear cost, without ever materializing a graph over the full dataset.

\noindent\textbf{From size balance to load balance.} The coloring objective enforces \emph{size} balance, but graph search induces a second, subtler imbalance: a small set of ``hub'' vectors are visited far more often than others during routing. If these hubs---and the traffic they attract---concentrate on a few shards, those shards become compute bottlenecks even though every shard holds an equal number of points. \sysname therefore treats \emph{query load} as a first-class \textit{capacity constraint} during L0 assignment: rather than sending every point to its top-voted shard, \sysname bounds the aggregate expected load each shard may absorb and spills excess points to their next-best shard. Hot regions are thus rebalanced before they can form a bottleneck, converting the residual load skew that survives size balancing into an explicitly controlled quantity.

\subsection{Data Multi-Assignment}
So far, $\sigma$ assigns each point to a single shard, which corresponds to an \emph{edge-cut} formulation. However, graph partitioning theory offers a powerful relaxation: the \emph{vertex-cut} (or overlapping) formulation. 

By allowing a point to be assigned to multiple shards ($\sigma:X\to 2^{\{1,\dots,P\}}$), we can replicate points that sit on heavily trafficked boundaries. This mathematically reduces the cut size, as queries no longer need to cross a partition boundary to reach a replicated neighbor. In \sysname, this translates to \emph{data multi-assignment}: selectively replicating boundary neighborhoods onto their next-nearest shards. This strictly lowers the required fan-out without sacrificing recall, at the cost of a bounded, user-configurable storage expansion.

With the offline index constructed and partitioned, the remaining challenge is how to efficiently route online queries against this topology. We detail \sysname's online navigational routing in the next section.

